%%%
%
% $Autor: C. Crety, H. Mey, J. Perewersenko $
% $Datum: 2025-07-10 $

% !TeX spellcheck = de_DE/GB
% !TeX program = pdflatex
% !BIB program = biber/bibtex
% !TeX encoding = utf8
%
%%%



\chapter{Setup für den Arduino Nano 33 BLE Sense}

\section{Einführung}

In diesem Kapitel wird der Prozess beschrieben, wie der Arduino Nano 33 BLE Sense mit einem Computer verbunden und die notwendigen Einstellungen vorgenommen werden, um mit der Programmierung zu beginnen.

Behandelt werden:
\begin{itemize}
	\item die Installation erforderlicher Tools
	\item die die Konfiguration in der Arduino IDE
	\item ein einfacher Beispiel-Sketch zur Überprüfung der Funktion
\end{itemize}

Zur Überprüfung der Konfiguration kann einer der vordefinierten Beispiel-Sketches genutzt werden, z.B. der einfache \FILE{blnik.ino}, der mit der IDE ausgeliefert wird.

\section{Konfiguration für den Arduino Nano 33 BLE Sense}

Um den \textbf{Arduino Nano 33 BLE Sense} in einer offline Umgebung zu programmieren sind folgende Schritte notwendig:

\begin{enumerate}
	\subsection{Installation des Treibers}
	\item \textbf{Arduino IDE installieren:} Installieren Sie die aktuelle Version der Arduino IDE auf Ihrem Computer.
	\item \textbf{Installation der Boardpakete:} Öffnen Sie die \menu{Arduino IDE} und wählen Sie im Menü \menu{Tools}, in der oberen linken Ecke.
	\item Im Menü \menu{Tools}, wählen Sie den \menu{Board Manager}.
	\item Suchen Sie im Fenster \menu{Board Manager} nach \textbf{Arduino Nano 33 BLE Sense}.
	\item Wählen Sie in den Suchergebnissen das Paket \menu{Arduino Mbed OS Nano Boards} und clicken Sie \menu{Install}, um das Paket zu installieren.
	\item Schließen Sie das Board über einen USB-A zu Micro-USB Kabel an den Computer an.
\end{enumerate}

\begin{center}
	\begin{tikzpicture}
		\node (G) at (0,0) {\includegraphics[width=8cm]{Arduino/ArdiunoIDE2/boardmanager1.png}};
		
		\draw[red,domain=0:360,line width=1.4pt] plot ({-2.6+1*cos(\x)},{0.7
			+0.5*sin(\x)});
		
	
		
		
	\end{tikzpicture}
				\captionof{figure}{Arduino Mbed OS Nano Board Installation}\label{fig:ArduinoMbedOSNanoBoardsInstallation}
\end{center}	



%%%%%%%%%%%%%%%%%%%%%%%%%%%%%%%%%%%%%%%%%%%%%%%%%%%%%%%%%%%%%%%%%%%%%%%%%%%%%%%%%%%%%%%

\subsection{Verbindung und Konfiguration des Arduino-Boards mit dem Computer}

%%%%%%%%%%%%%%%%%%%%%%%%%%%%%%%%%%%%%%%%%%%%%%%%%%%%%%%%%%%%%%%%%%%%%%%%%%%%%%%%%%%%%%%

Um entweder ein vordefiniertes Beispiel oder ein eigenes Programm auf dem Arduino Nano 33 BLE Sense auszuführen, befolgen Sie die folgenden Schritte::

\begin{enumerate}
	
	\item Verbinden Sie das Arduino-Board mit dem Computer.
	\item Öffnen Sie die Arduino IDE, dann erscheint ein leeres Projekt mir den Standardfunktionen \PYTHON{void setup()} und \PYTHON
	{void loop()}.
	\item Gehen Sie im Menü zu \menu{Tools}, wählen Sie \menu{Board} und wählen Sie das angeschlossene Board \textbf{Arduino Nano 33 BLE Sense}, wie in Abbildung~\ref{fig:SelecttheConnectedboard-hereArduinoNano33BLESense}.
	
\end{enumerate}

Damit ist die IDE so konfiguriert, dass sie das Board erkennt und mit ihm kommunizieren kann.

\begin{figure}[H]\centering
	\includegraphics[width=8cm]{Arduino/ArdiunoIDE2/SelectTheConnectedBoardArduinoNano33BLESense}
	\caption{Verbundenes Board auswählen - hier Arduino Nano 33 BLE Sense}
	\label{fig:SelecttheConnectedboard-hereArduinoNano33BLESense}		
\end{figure}

%%%%%%%%%%%%%%%%%%%%%%%%%%%%%%%%%%%%%%%%%%%%%%%%%%%%%%%%%%%%%%%%%%%%%%%%%%%%%%%%%%%%%%%

\subsection{Auswahl des passenden Ports}

%%%%%%%%%%%%%%%%%%%%%%%%%%%%%%%%%%%%%%%%%%%%%%%%%%%%%%%%%%%%%%%%%%%%%%%%%%%%%%%%%%%%%%%

Nach der Auswahl des Arduino Nano 33 BLE Sense, muss nun der richtige Port ausgewählt werden. Dafür versetzen Sie das Board in den Bootloader-Modus, drücken Sie dazu den weißen Reset-Knopf (siehe Abbildung ~\ref{fig:ArduinoNano33BLESenseResetButton}).

\begin{center}
	\begin{tikzpicture}
		\node (G) at (0,0) {\includegraphics[width=8cm]{Arduino/ArdiunoIDE2/resetbutton.png}};
		
		\draw[green,domain=0:360,line width=1.4pt] plot ({-1.5+0.4*cos(\x)},{-1.6
			+0.4*sin(\x)});
	\end{tikzpicture}
		\captionof{figure}{Arduino Nano 33 BLE Sense Reset Button}\label{fig:ArduinoNano33BLESenseResetButton}		
\end{center}


Wenn der Bootloader-Modus aktiviert ist, leuchtet die orangefarbene LED, wie in Abbildung ~\ref{fig:ArduinoNano33BLESenseResetButton}			



\begin{center}
	\begin{tikzpicture}
		\node (G) at (0,0) {\includegraphics[width=8cm]{Arduino/ArdiunoIDE2/resetbutton.png}};
		
		\draw[green,domain=0:360,line width=1.4pt] plot ({-2.8+0.4*cos(\x)},{-0.8
			+0.4*sin(\x)});
			
		\draw[orange,domain=0:360,line width=1.4pt] plot ({-2.8+0.4*cos(\x)},{-2.6
	+0.4*sin(\x)});
			
			
			
	\end{tikzpicture}
		\captionof{figure}{grüne und orangene Builtin-LED}\label{fig:ArduinoNano33BLESenseResetButton}	
\end{center}	


Nun müssen Sie den Port auswählen. Gehen Sie dazu zu \menu{Tools}, wählen Sie \menu{Port} und versichern Sie sich, dass der richtig Port ausgewählt ist. ~\ref{fig:SelectAvailablePortforUploadingArduinoSketch}	



\begin{figure}[H]\centering
	\includegraphics[width=8cm]{Arduino/ArdiunoIDE2/SelectAvailablePortForUploadingArduinoSketch}
	\caption{Zum Upload von Arduino Sketches verfügbaren Port auswählen}
	\label{fig:SelectAvailablePortforUploadingArduinoSketch}		
\end{figure}




\subsection{Sketch hochladen und überprüfen}\label{uploadcode}

Nachdem der passende Port ausgewählt wurde, folgt nun der nächste Schritt: das Hochladen eines Sketches.
Vor dem Hochladen sollte der Sketch mit der Schaltfläche "Verify" kompiliert werde, dadurch wird geprüft, ob der Code fehlerfrei ist. Nach erfolgreicher Überprüfung kann der Sketch mit der Schaltfläche "Upload" an das Arduino Board übertragen werden.(Abbildung ~\ref{fig:Upload the Program in Arduino board})

Als nächstes muss die Verbingung bestätigt werden. Unten rechts im Fenster der Arduino IDE sollte der Name des angeschlossenen Boards und der verwendete Port angezeigt werden, das bestätigt, dass das Board korrekt erkannt wurde.

\begin{figure}[H]\centering
	\includegraphics[width=8cm]{Arduino/ArdiunoIDE2/UploadTheProgramInArduinoBoard}
	\caption{Upload des Sketches auf das Arduino Board}
	\label{fig:Upload the Program in Arduino board}		
\end{figure}

Nach dem Hochladen wird der Sketch kompiliert. Falls dabei Fehler auftreten, werden sie im unteren schwarzen Konsolenbereich der Arduino IDE angezeigt. Nach dem erfolgreichen Upload sollte der verwendete Port erneut ausgewählt werden, um z.B. die Ausgabe im seriellen Monitor anzeigen zu können.

\begin{figure}[H]\centering
	\includegraphics[width=8cm]{Arduino/ArdiunoIDE2/SettingThePort}
	\caption{Auswählen des Ports}
	\label{fig:SettingthePort}		
\end{figure} 

\section{Konfiguration testen}

In diesem Abschnitt wird überprüft, ob das BOard korrekt mit der Arduino IDE verbunden ist. Dazu wird ein einfacher Beispiel-Sketch ("Blink") ausgeführt, bei dem eine eingebaute LED blinkt. So lässt sich sowohl die Verbindung als auch die Grundfunktionalität des Boards überprüfen.

\subsection{Test anhand des Beispiel-Sketches \FILE{blink.ino}}

Die Arduino-IDE enthält eine Reihe integrierter Beispiele zu Testzwecken.
Um die Konfiguration zu überprüfen und das Board einzurichten, kann zunächst das einfache Beispiel \FILE{blink.ino} geöffnet werden.
Nach dem Hochladen des Beispiels \FILE{blink.ino} blinkt die eingebaute LED mit 2 Hz.

Dafür werden folgende Schritte ausgeführt:

\begin{enumerate}
\item \textbf{Beispiele öffnen:} In der Arduino IDE auf \menu{File} \menu{Examples} klicken. 

\item \textbf{Sketch auswählen:} Eine Liste von Beispiel-Sketches wird angezeigt. Im Untermenü \menu{01.Basics} \menu{Blink} auswählen, wie in Abbildung ~\ref{fig:LED-ExampleTest} gezeigt.

\item \textbf{Blink Beispiel-Sketch:} Der Beispiel-Sketch \FILE{blink.ino} erscheint in einem neuen Fenster der Arduino IDE ~\ref{fig:LED-Built_Code}.

\item \textbf{Sketch hochladen:} Mit einem Klick auf den \menu{Upload} Button wird der Sketch hochgeladen. Wenn der Upload erfolgreich ist, dann beginnt die eingebaute LED mit einer Frequenz von 2 Hz zu blinken. 
\end{enumerate}


\begin{center}
	\includegraphics[width=8cm]{Arduino/ArdiunoIDE2/blink_example_IDE2.3.3.png}
	\captionof{figure}{LED Beispiel-Sketch}\label{fig:LED-ExampleTest}		
\end{center} 



\begin{center}
	\includegraphics[width=8cm]{Arduino/ArdiunoIDE2/ledblink_code.png}
	\captionof{figure}{LED-Built Example}\label{fig:LED-Built_Code}		
\end{center} 

Wenn der Test erfolgreich ist, dann ist somit bestätigt, dass die Verbindung zur Arduino IDE funktioniert und die Hardware betriebsbereit ist.

